% ============ 8 层的直像与逆像 (原书 pp.58-63) ============
\chapter{层的直像与逆像}\label{sec:8}
\index{直像}
\index{逆像}
\index{零延拓}
\index{内射对象}
\index{支集}
\index{紧支集}

\subsection*{层的直像}

设 $\varphi \colon Y \to X$ 是簇(或概形)的态射，$P$ 是 $Y_{\mathrm{et}}$ 上的预层.对平展映射 $U \to X$，定义
\[
\varphi_{*}P(U) = P(U \times_{X} Y).
\]
由于 $U \times_{X} Y \to Y$ 平展(见命题~\ref{prop:2.11})，这个定义有意义.配上显然的限制映射，$\varphi_{*}P$ 成为 $X_{\mathrm{et}}$ 上的预层.

\begin{lemma}\label{lem:8.1}
若 $F$ 是层，则 $\varphi_{*}F$ 也是层.
\end{lemma}

\begin{proof}
对 $X$ 上的簇(或概形) $V$，设 $V_Y$ 表示 $Y$ 上的簇(或概形) $V \times_{X} Y$.则 $V \mapsto V_Y$ 是把平展映射送到平展映射，把满映射族送到满映射族，把 $X$ 上的纤维积送到 $Y$ 上的纤维积的函子.

设 $(U_i \to U)$ 是 $\mathbf{Et}/X$ 中平展映射的满族.则 $(U_{iY} \to U_Y)$ 是 $\mathbf{Et}/Y$ 中平展映射的满族，从而
\[
F(U_Y) \to \prod F(U_{iY}) \rightrightarrows \prod F(U_{iY} \times_{Y} U_{jY})
\]
正合.但这等于序列
\[
(\varphi_{*}F)(U) \to \prod (\varphi_{*}F)(U_i) \rightrightarrows \prod (\varphi_{*}F)(U_i \times_{X} U_j),
\]
因此它也正合.
\end{proof}

显然，函子
\[
\varphi_{*} \colon \mathrm{PreSh}(Y_{\mathrm{et}}) \to \mathrm{PreSh}(X_{\mathrm{et}})
\]
正合.因此，它的限制
\[
\varphi_{*} \colon \mathrm{Sh}(Y_{\mathrm{et}}) \to \mathrm{Sh}(X_{\mathrm{et}})
\]
左正合.它通常不右正合：$F \to F'$ 局部满不蕴含 $\varphi_{*}F \to \varphi_{*}F'$ 局部满.例如，若 $X$ 是代数闭域上的代数簇，$\varphi \colon X \to P$ 是把 $X$ 映到点的映射，则 $\varphi_{*}$ 本质上就是把层 $F$ 送到其整体截面群 $\Gamma(X, F)$ 的函子，而 $F \to F'$ 局部满不蕴含 $\Gamma(X, F) \to \Gamma(X, F')$ 满射.

\begin{example}\label{ex:8.2}
设 $i \colon \bar{x} \to X$ 是 $X$ 的几何点.函子 $F \mapsto F(\bar{x})$ 把 $\bar{x}$ 上的层范畴等同于 Abel 群范畴.设 $\Lambda$ 是看作 $\bar{x}$ 上层的 Abel 群.则 $i_{*}\Lambda = \Lambda_{\bar{x}}$，即 \S 6 中定义的摩天楼层.

$Y$ 的几何点 $\bar{y} \to Y$ 通过 $\bar{y} \to Y \xrightarrow{\varphi} X$ 定义 $X$ 的几何点，记为 $\bar{x}$(或 $\varphi(\bar{y})$).显然
\[
(\varphi_{*}F)_{\bar{x}} = \varinjlim F(V),
\]
其中极限遍历 $\bar{y}$ 的所有这样的平展邻域：对 $\bar{x}$ 的某个平展邻域 $U$，它们具有形式 $U_Y$.于是，有典范映射
\[
(\varphi_{*}F)_{\bar{x}} \to F_{\bar{y}}.
\]
一般地，这个映射既不是单射也不是满射.
\end{example}

\begin{proposition}\label{prop:8.3}
(a) 设 $\varphi \colon V \hookrightarrow X$ 是开浸入，即开子簇(或子概形)到 $X$ 的包含.则
\[
(\varphi_{*}F)_{\bar{x}} = \begin{cases}
F_{\bar{x}} & x \in V;\\
? & x \notin V.
\end{cases}
\]
(b) 设 $\varphi \colon Z \hookrightarrow X$ 是闭浸入，即闭子簇(或子概形)到 $X$ 的包含.则
\[
(\varphi_{*}F)_{\bar{x}} = \begin{cases}
F_{\bar{x}} & x \in Z;\\
0 & x \notin Z.
\end{cases}
\]
(c) 设 $\varphi \colon Y \to X$ 是有限映射.则
\[
(\varphi_{*}F)_{\bar{x}} = \bigoplus_{y \mapsto x} F^{d(y)}_{\bar{y}},
\]
其中 $d(y)$ 是 $\kappa(y)$ 在 $\kappa(x)$ 上的可分次数.例如，若 $\varphi$ 是代数闭域上次数为 $d$ 的簇的有限平展映射，则
\[
(\varphi_{*}F)_{\bar{x}} = F^{d}_{\bar{x}}.
\]
\end{proposition}

\begin{proof}
当 $\varphi \colon V \hookrightarrow X$ 是浸入(开的或闭的)时，对任何态射 $\varphi' \colon U \to X$，纤维积 $U \times_{X} V = \varphi'^{-1}(V)$.

(a) 若 $x \in V$，则对 $\bar{x}$ 的任何``足够小''的平展邻域 $\varphi' \colon U \to X$, $\varphi'(U) \subset V$，从而 $U = \varphi'^{-1}(V) = U \times_{X} V$.因此，$\bar{x}$ 的具有形式 $U_V$ 的平展邻域构成共尾集，这证明了第一个等式.关于第二个，除了指出当 $x \notin V$ 时 $(\varphi_{*}F)_{\bar{x}}$ 不必为零(见下面的例子)外，没什么可说的.

(b) 若 $x \notin Z$，则对 $\bar{x}$ 的任何``足够小''的平展邻域 $\varphi' \colon U \to X$, $\varphi'(U) \cap Z = \varnothing$，从而 $U \times_{X} Z = \varphi'^{-1}(Z) = \varnothing$；因此 $F(U_Z) = 0$.当 $x \in Z$ 时，必须看到每个平展映射 $\bar{\varphi}' \colon \bar{U} \to Z$($x \in \varphi'(\bar{U})$)都``延拓''为平展映射 $\varphi' \colon U \to X$.用环的语言，这相当于证明平展同态 $\bar{A} \to \bar{B}$($\bar{A} = A/\mathfrak{a}$)提升为平展同态 $A \to B$.但这很容易：我们可以假定 $\bar{B} = (\bar{A}[T]/(\bar{f}(T)))_{\bar{b}}$，其中 $\bar{f}(T)$ 是首一多项式，使得 $\bar{f}'(T)$ 在 $(\bar{A}[T]/(\bar{f}(T)))_{\bar{b}}$ 中可逆---见 \S 2(标准平展态射)；选取提升 $\bar{f}(T)$ 的 $f(T) \in A[T]$，对适当的 $b$ 置 $B = (A[T]/(f(T)))_b$.

(c) 我们省略证明.
\end{proof}

\begin{corollary}\label{cor:8.4}
若 $\varphi$ 有限或是闭浸入，则函子 $\varphi_{*}$ 正合.
\end{corollary}

\begin{proof}
这由该命题与命题~\ref{prop:7.6} 推出.
\end{proof}

\begin{example}\label{ex:8.5}
我们列出一些例子，都是类比.

(a) 设 $X$ 是以原点为中心的 $\mathbb{R}^{2}$ 中半径 1 的开圆盘，$U = X \setminus \{(0,0)\}$，$j$ 是包含 $U \hookrightarrow X$.设 $F$ 是对应于 $\pi_1(U, u)$-模 $F$ 的 $U$ 上的局部常值层---回忆 $\pi_1(U, u) = \mathbb{Z}$(见命题~\ref{prop:6.15}；这里 $u$ 是 $U$ 的任意点).则
\[
(j_{*}F)_{(0,0)} = F^{\pi_1(U,u)} \quad (\text{被 } \pi_1 \text{ 固定的元素}).
\]

(b) 设 $X$ 是代数闭域上的仿射线，$U = X \setminus \{0\}$，$j$ 是包含 $U \hookrightarrow X$.设 $F$ 是对应于 $\pi_1(U, \bar{u})$-模 $F$ 的 $U$ 上的局部常值层---回忆 $\pi_1(U, \bar{u}) = \widehat{\mathbb{Z}}$.则
\[
(j_{*}F)_{\bar{0}} = F^{\pi_1(U,\bar{u})} \quad (\text{被 } \pi_1 \text{ 固定的元素}).
\]

(c) 设 $X = \operatorname{Spec} R$，其中 $R$ 是 Hensel 离散赋值环，$U = \operatorname{Spec} K$，其中 $K$ 是 $R$ 的分式域，$j$ 是包含 $U \hookrightarrow X$.设 $F$ 是对应于 $\pi_1(U, \bar{u})$-模 $F$ 的 $U$ 上的局部常值层---回忆 $\pi_1(U, \bar{u}) = \operatorname{Gal}(K^{\mathrm{sep}}/K)$.这里 $\bar{u}$ 是对应于包含 $K \hookrightarrow K^{\mathrm{sep}}$ 的几何点.设 $I$(惯性群)是 $\pi_1(U, \bar{u})$ 中平凡作用于剩余域的元素组成的子群.则
\[
(j_{*}F)_{\bar{x}} = F^{I} \quad (\text{被 } I \text{ 固定的元素}),
\]
其中 $x$ 是 $\operatorname{Spec} R$ 的闭点.
\end{example}

\begin{proposition}\label{prop:8.6}
对任意态射 $Z \xrightarrow{\varphi} Y \xrightarrow{\varphi'} X$，$(\varphi' \circ \varphi)_{*} = \varphi'_{*} \circ \varphi_{*}$.
\end{proposition}

\begin{proof}
这由定义显然.
\end{proof}

\subsection*{层的逆像}

设 $\varphi \colon Y \to X$ 是簇(或概形)的态射.我们将定义函子 $\varphi_{*}$ 的左伴随.设 $P$ 是 $X_{\mathrm{et}}$ 上的预层.对平展映射 $V \to Y$，定义
\[
P'(V) = \varinjlim P(U),
\]
其中正向极限遍历交换图
\[
\begin{tikzcd}[row sep=small, column sep=small]
  V \arrow[r] \arrow[rd] & U \arrow[d] \\
                         & X
\end{tikzcd}
\]
$U \to X$ 平展，且映射 $V \to U$ 与 $V \to Y$ 使图交换.容易看出，对 $Y$ 上的任意预层 $Q$，下列对象之间有自然的双射对应：
\begin{itemize}
  \item 态射 $P' \to Q$；
  \item 由如上交换图指标化的映射族 $P(U) \to Q(V)$，与限制映射相容；
  \item 态射 $P \to \varphi_{*}Q$.
\end{itemize}
因此
\[
\operatorname{Hom}_{Y_{\mathrm{et}}}(P', Q) \cong \operatorname{Hom}_{X_{\mathrm{et}}}(P, \varphi_{*}Q),
\]
关于 $P$ 与 $Q$ 函子性地成立.遗憾的是，即使 $P$ 是层，$P'$ 也不必是层.因此，对 $X_{\mathrm{et}}$ 上的层 $F$，定义 $\varphi^{*}F = a(F')$.则对 $Y_{\mathrm{et}}$ 上的任意层 $G$，
\[
\operatorname{Hom}_{Y_{\mathrm{et}}}(\varphi^{*}F, G) \cong \operatorname{Hom}_{Y_{\mathrm{et}}}(F', G) \cong \operatorname{Hom}_{X_{\mathrm{et}}}(F, \varphi_{*}G),
\]
从而 $\varphi^{*}$ 是 $\varphi_{*} \colon \mathrm{Sh}(Y_{\mathrm{et}}) \to \mathrm{Sh}(X_{\mathrm{et}})$ 的左伴随.

\begin{proposition}\label{prop:8.7}
对任意态射 $Z \xrightarrow{\varphi} Y \xrightarrow{\varphi'} X$，$(\varphi' \circ \varphi)^{*} = \varphi^{*} \circ \varphi'^{*}$.
\end{proposition}

\begin{proof}
两者都是 $(\varphi' \circ \varphi)_{*} = \varphi'_{*} \circ \varphi_{*}$ 的左伴随.
\end{proof}

\begin{example}\label{ex:8.8}
设 $\varphi \colon U \to X$ 是平展态射.对 $X_{\mathrm{et}}$ 上的任意层 $F$ 与 $U_{\mathrm{et}}$ 上的任意层 $G$，容易看出
\[
\operatorname{Hom}(F|_{U_{\mathrm{et}}}, G) \cong \operatorname{Hom}(F, \varphi_{*}G)
\]
---在两种情形中，给出态射都是给出由平展映射 $V \to U$ 指标化的，与限制相容的映射族 $F(V) \to G(V)$.因此，在这种情况下，$\varphi^{*} \colon \mathrm{Sh}(X_{\mathrm{et}}) \to \mathrm{Sh}(U_{\mathrm{et}})$ 就是限制映射.

这也可以直接从 $\varphi^{*}$ 的定义看出，因为对每个 $V \to U$，取极限的图族中有终对象，即
\[
\begin{tikzcd}[row sep=small, column sep=small]
  V \arrow[r, equal] \arrow[rd] & V \arrow[d] \\
  U \arrow[r, "\varphi"] & X
\end{tikzcd}
\]
\end{example}

\begin{remark}\label{rem:8.9}
设 $i \colon \bar{x} \to X$ 是 $X$ 的几何点.对 $X_{\mathrm{et}}$ 上的任意层 $F$，$(i^{*}F)(\bar{x}) = F_{\bar{x}}$---这由 $i^{*}$ 与 $F_{\bar{x}}$ 的定义显然.因此，对任意态射 $\varphi \colon Y \to X$ 与 $Y$ 的几何点 $i \colon \bar{y} \to Y$，
\[
(\varphi^{*}F)_{\bar{y}} = i^{*}(\varphi^{*}F)(\bar{y}) = F_{\bar{x}},
\]
其中 $\bar{x}$ 是 $X$ 的几何点 $\bar{y} \xrightarrow{i} Y \xrightarrow{\varphi} X$.

由于这对 $Y$ 的所有几何点都成立，我们看到 $\varphi^{*}$ 正合，从而 $\varphi_{*}$ 保持内射对象(由 7.3).
\end{remark}

\begin{remark}\label{rem:8.10}
设 $F_Z$ 是由簇 $Z$ 定义的 $X_{\mathrm{et}}$ 上的层.则 $\varphi^{*}(F_Z)$ 是由簇 $Z \times_{Y} X$ 定义的层 $F_{Z \times_{Y} X}$ 这一点不总为真.例如，若 $\varphi \colon X \to P$ 是代数闭域 $k$ 上的簇到点的映射，则 $\varphi^{*}(\mathbb{G}_{a})$ 是由群 $\mathbb{G}_{a}(k) = k$ 定义的常值层.

然而，当 $\varphi$ 平展或 $Z \to X$ 平展时，$\varphi^{*}(F_Z) = F_{Z \times_{Y} X}$ 为真.
\end{remark}

\begin{remark}\label{rem:8.11}
对拓扑空间的连续映射 $\varphi \colon Y \to X$，可以用相伴的 \'etal\'e 空间非常简单地描述层的逆像：$\varphi^{*}(F)$ 是如下层：它在开子集 $U \subset Y$ 上的截面是连续映射 $s \colon U \to E(F)$，使得对所有 $u$，$s(u)$ 位于 $\varphi(u)$ 上的茎中.
\end{remark}

\subsection*{足够内射对象的存在性}

设 $X$ 是簇(或概形).

\begin{proposition}\label{prop:8.12}
$\mathrm{Sh}(X_{\mathrm{et}})$ 中的每个层 $F$ 都能嵌入内射层.
\end{proposition}

\begin{proof}
对每个 $x \in X$，选取像为 $x$ 的几何点 $i_x \colon \bar{x} \to X$，并把 Abel 群 $F_{\bar{x}}$ 嵌入内射 Abel 群 $I(x)$.则 $I_x \overset{\mathrm{def}}{=} i_{x*}(I(x))$ 内射(见注记~\ref{rem:8.9}).由于内射对象的乘积内射，$I \overset{\mathrm{def}}{=} \prod I_x$ 是内射层.复合 $F \hookrightarrow P^{*} \hookrightarrow I$ 就是所求的嵌入.
\end{proof}

\subsection*{零延拓}

设 $X$ 是簇(或概形)，$j \colon U \hookrightarrow X$ 是开浸入.如我们上面所述，对 $U_{\mathrm{et}}$ 上的层 $F$，$j_{*}F$ 在 $U$ 之外的点处的茎不必为零.现在我们定义``零延拓''函子 $j_{!}$，使得 $j_{!}F$ 确实有这个性质.

设 $P$ 是 $U_{\mathrm{et}}$ 上的预层.对任何平展映射 $\varphi' \colon V \to X$，定义
\[
P_{!}(V) = \begin{cases}
P(V) & \text{若 } \varphi'(V) \subset U;\\
0 & \text{其他情形}.
\end{cases}
\]
则 $P_{!}$ 是 $X_{\mathrm{et}}$ 上的预层，且对 $X_{\mathrm{et}}$ 上的任何其他预层 $Q$，态射 $P \to Q|_{U}$ 唯一延拓为态射 $P_{!} \to Q$(显然).因此
\[
\operatorname{Hom}_{X_{\mathrm{et}}}(P_{!}, Q) \cong \operatorname{Hom}_{U_{\mathrm{et}}}(P, Q|_{U})
\]
函子性地成立.遗憾的是，即使 $P$ 是层，$P_{!}$ 也不必是层.因此，对 $U_{\mathrm{et}}$ 上的层 $F$，定义 $j_{!}F = a(F_{!})$.则对 $X_{\mathrm{et}}$ 上的任意层 $G$，
\[
\operatorname{Hom}_{X_{\mathrm{et}}}(j_{!}F, G) \cong \operatorname{Hom}_{X_{\mathrm{et}}}(F_{!}, G) \cong \operatorname{Hom}_{U_{\mathrm{et}}}(F, G|_{U}),
\]
从而 $j_{!}$ 是 $j^{*} \colon \mathrm{Sh}(X_{\mathrm{et}}) \to \mathrm{Sh}(U_{\mathrm{et}})$ 的左伴随.

\begin{proposition}\label{prop:8.13}
设 $j \colon U \hookrightarrow X$ 是开浸入.对 $U_{\mathrm{et}}$ 上的任意层 $F$ 与几何点 $\bar{x} \to X$,
\[
(j_{!}F)_{\bar{x}} = \begin{cases}
F_{\bar{x}} & x \in U;\\
0 & x \notin U.
\end{cases}
\]
\end{proposition}

\begin{proof}
由引理~\ref{lem:7.14}，只需对把 $j_{!}F$ 换成 $F_{!}$ 证明，这时它立即由定义推出.
\end{proof}

\begin{corollary}\label{cor:8.14}
函子 $j_{!} \colon \mathrm{Sh}(U_{\mathrm{et}}) \to \mathrm{Sh}(X_{\mathrm{et}})$ 正合，且 $j^{*}$ 保持内射对象.
\end{corollary}

\begin{proof}
论断的第一部分由该命题与命题~\ref{prop:7.6} 推出，第二部分由第一部分与~\ref{item:7.3} 推出.
\end{proof}

设 $Z$ 是 $U$ 在 $X$ 中的补集，用 $i$ 表示包含 $Z \hookrightarrow X$.设 $F$ 是 $X_{\mathrm{et}}$ 上的层.有典范态射 $j_{!}j^{*}F \to F$(由伴随性对应于 $j^{*}F$ 上的恒等映射)与典范态射 $F \to i_{*}i^{*}F$(由伴随性对应于 $i^{*}F$ 上的恒等映射).

\begin{proposition}\label{prop:8.15}
对 $X$ 上的任意层 $F$，序列
\[
0 \to j_{!}j^{*}F \to F \to i_{*}i^{*}F \to 0
\]
正合.
\end{proposition}

\begin{proof}
这可以在茎上检验.对 $x \in U$，茎序列是
\[
0 \to F_{\bar{x}} \xrightarrow{\mathrm{id}} F_{\bar{x}} \to 0 \to 0;
\]
对 $x \notin U$，茎序列是
\[
0 \to 0 \to F_{\bar{x}} \xrightarrow{\mathrm{id}} F_{\bar{x}} \to 0.
\]
两者显然正合.
\end{proof}

\begin{remark}\label{rem:8.16}
可以对任何平展映射 $j \colon U \to X$ 定义 $j_{!}$.设 $F$ 是 $U_{\mathrm{et}}$ 上的层.对任何平展映射 $\varphi' \colon V \to X$，定义
\[
F_{!}(V) = \bigoplus_{\alpha} F(V),
\]
其中求和遍历满足 $j \circ \alpha = \varphi'$ 的态射 $\alpha \colon V \to U$.则 $F_{!}$ 是 $X_{\mathrm{et}}$ 上的预层，我们定义 $j_{!}F$ 为它的相伴层.同样，$j_{!}$ 是 $j^{*}$ 的左伴随且正合.因此 $j^{*}$ 保持内射对象.
\end{remark}

\subsection*{\texorpdfstring{$X = U \cup Z$ 上的层}{X = U ∪ Z 上的层}}

设 $X$ 是簇(或概形)，
\[
U \xrightarrow{j} X \xleftarrow{i} Z
\]
是态射，其中 $j$ 是开浸入，$i$ 是闭浸入，且 $X = j(U) \cup i(Z)$.我们把 $U$ 与 $Z$ 等同于 $X$ 的子簇(或子概形).

由 $X_{\mathrm{et}}$ 上的层 $F$，得到 $Z$ 上的层 $F_1 \overset{\mathrm{def}}{=} i^{*}F$ 与 $U$ 上的层 $F_2 = j^{*}F$.此外，有典范同态 $F \to j_{*}j^{*}F$(由伴随性对应于 $j^{*}F$ 上的恒等映射).对它应用 $i^{*}$，得到态射 $\Phi_F \colon F_1 \to i^{*}j_{*}F_2$.

设 $\mathrm{Tr}(X, U, Z)$ 是如下范畴：对象是三元组 $(F_1, F_2, \Phi)$，其中 $F_1$ 是 $Z_{\mathrm{et}}$ 上的层，$F_2$ 是 $U_{\mathrm{et}}$ 上的层，$\Phi$ 是映射 $F_1 \to i^{*}j_{*}F_2$.态射 $(F_1, F_2, \Phi) \to (F'_1, F'_2, \Phi')$ 是满足下图交换的态射对 $\psi_1 \colon F_1 \to F'_1$ 与 $\psi_2 \colon F_2 \to F'_2$:
\[
\begin{tikzcd}
F_1 \arrow[r, "\Phi"] \arrow[d, "\psi_1"] & i^{*}j_{*}F_2 \arrow[d, "i^{*}j_{*}\psi_2"] \\
F'_1 \arrow[r, "\Phi'"'] & i^{*}j_{*}F'_2
\end{tikzcd}
\]

\begin{proposition}\label{prop:8.17}
函子 $F \mapsto (F_1, F_2, \Phi_F)$ 是等价 $\mathrm{Sh}(X_{\mathrm{et}}) \to \mathrm{Tr}(X, U, Z)$.
\end{proposition}

\begin{proof}
本质逆由函子
\[
(F_1, F_2, \Phi) \mapsto i_{*}F_1 \times_{i_{*}i^{*}j_{*}F_2} j_{*}F_2
\]
给出.细节见 EC p.74.
\end{proof}

在这个范畴等价下，
\[
0 \to j_{!}j^{*}F \to F \to i_{*}i^{*}F \to 0
\]
对应于正合序列
\[
0 \to (0, j^{*}F, 0) \to (i^{*}F, j^{*}F, \Phi_F) \to (i^{*}F, 0, 0) \to 0.
\]
当且仅当 $\Phi_F = 0$ 时该序列分裂.

设 $Y$ 是 $X$ 的子集.我们说层 $F$ 的支集在 $Y$ 中，如果只要 $x \notin Y$ 就有 $F_{\bar{x}} = 0$.

\begin{corollary}\label{cor:8.18}
对任何闭浸入 $i \colon Z \hookrightarrow X$，$i_{*}$ 定义了 $Z_{\mathrm{et}}$ 上层范畴与支集在 $Z$ 中的 $X_{\mathrm{et}}$ 上层范畴之间的等价.
\end{corollary}

\begin{proof}
这由该命题与显然的事实推出：$F$ 的支集在 $Z$ 中，当且仅当它对应于形如 $(F_1, 0, 0)$ 的三元组.
\end{proof}
